\begin{thebibliography}{99}
\bibitem{bib-bib1} 
\par\noindent    R. F. Blute, J. R. B. Cockett, J.-S. P. Lemay, and R. A. G. Seely.  Differential categories revisited.  Applied Categorical Structures, 2020.  doi:10.1007/s10485-019-09572-y.  

\bibitem{bib-bib2} 
\par\noindent    R. F. Blute, J. R. B. Cockett, and R. A. G. Seely.  Differential categories.  Mathematical Structures in Computer Science, 2006.  doi:10.1017/S0960129506005676.  

\bibitem{bib-bib3} 
\par\noindent    R. F. Blute, J. R. B. Cockett, and R. A. G. Seely.  Cartesian differential categories.  Theory and Applications of Categories, 2009.  

\bibitem{bib-bib4} 
\par\noindent    Richard Blute, Thomas Ehrhard, and Christine Tasson.  A convenient differential category.  Cahiers de topologie et géométrie différentielle catégoriques, 53(3):211–232, 2012.  arXiv:1006.3140.  

\bibitem{bib-bib5} 
\par\noindent    A. Brunel, D. Mazza, and M. Pagani.  Backpropagation in the Simply Typed Lambda-calculus with Linear Negation.  POPL, 2020.  doi:10.1145/3371132.  

\bibitem{bib-bib6} 
\par\noindent    Antonio Bucciarelli, Thomas Ehrhard, and Giulio Manzonetto.  Not enough points is enough.  In Computer Science Logic, 2007.  doi:10.1007/978-3-540-74915-8\_24.  

\bibitem{bib-bib7} 
\par\noindent    J. R. B. Cockett and J.-S. Lemay.  Integral categories and calculus categories.  Mathematical Structures in Computer Science, 29:243–308, 2019.  doi:10.1017/S0960129518000014.  

\bibitem{bib-bib8} 
\par\noindent    J. Robin B. Cockett and Geoff S. H. Cruttwell.  Differential Structure, Tangent Structure, and SDG.  Appl. Categorical Struct., 22(2):331–417, 2014.  doi:10.1007/s10485-013-9312-0.  

\bibitem{bib-bib9} 
\par\noindent    J. Robin B. Cockett, Geoff S. H. Cruttwell, Jonathan Gallagher, Jean-Simon Pacaud Lemay, Benjamin MacAdam, Gordon D. Plotkin, and Dorette Pronk.  Reverse derivative categories.  In Maribel Fernández and Anca Muscholl, editors, 28th EACSL Annual Conference on Computer Science Logic, CSL 2020, January 13-16, 2020, Barcelona, Spain, volume 152 of LIPIcs, pages 18:1–18:16. Schloss Dagstuhl – Leibniz-Zentrum für Informatik, 2020.  doi:10.4230/LIPIcs.CSL.2020.18.  

\bibitem{bib-bib10} 
\par\noindent    Y. Dabrowski and M. Kerjean.  Models of Linear Logic based on the Schwartz $\epsilon$ product.  Theory and Applications of Categories, 2020.  doi:10.48550/arXiv.1712.07344.  

\bibitem{bib-bib11} 
\par\noindent    T. Ehrhard.  An introduction to differential linear logic: proof-nets, models and antiderivatives.  Math. Struct. Comput. Sci., 2017.  doi:10.1017/S0960129516000372.  

\bibitem{bib-bib12} 
\par\noindent    T. Ehrhard and L. Regnier.  Differential interaction nets.  Theoretical Computer Science, 364(2), 2006.  doi:10.1016/j.tcs.2006.08.003.  

\bibitem{bib-bib13} 
\par\noindent    Conal Elliott.  The simple essence of automatic differentiation.  In Proceedings of the ACM on Programming Languages (ICFP), 2018.  doi:10.1145/3236765.  

\bibitem{bib-bib14} 
\par\noindent    M. P. Fiore.  Differential structure in models of multiplicative biadditive intuitionistic linear logic.  In TLCA, 2007.  doi:10.1007/978-3-540-73228-0\_13.  

\bibitem{bib-bib15} 
\par\noindent    J.-Y. Girard.  Linear logic.  Theoret. Comput. Sci., 50(1), 1987.  doi:10.1016/0304-3975(87)90045-4.  

\bibitem{bib-bib16} 
\par\noindent    Marie Kerjean.  A logical account for linear partial differential equations.  In Anuj Dawar and Erich Grädel, editors, Proceedings of the 33rd Annual ACM/IEEE Symposium on Logic in Computer Science, LICS 2018, Oxford, UK, July 09-12, 2018, pages 589–598. ACM, 2018.  doi:10.1145/3209108.3209192.  

\bibitem{bib-bib17} 
\par\noindent    Paul-André Mellies.  Categorical semantics of linear logic.  Société Mathématique de France, 2008.  

\bibitem{bib-bib18} 
\par\noindent    L. Schwartz.  Théorie des distributions.  Hermann, Paris, 1966.  

\bibitem{bib-bib19} 
\par\noindent    R. A. G. Seely.  Linear Logic, *-Autonomous Categories and Cofree Coalgebras.  In In Categories in Computer Science and Logic. American Mathematical Society, 1989.  

\end{thebibliography}
